The scoreboard's invariance to $B$ also clarifies how absence clamping interacts with the alert threshold. Once never-mentioned factors are fixed at zero, the residual free energy of the BN is concentrated on a small parent set. Additional inspections within that set update $p$ continuously, which is visible in the mild Brier improvement from $B{=}1$ to $B{\geq}2$, but they rarely flip the discrete alert decision once $\tau{=}0.51$ has been crossed or failed. Classification metrics can therefore plateau while probabilistic scores continue to move---a distinction that accuracy-only reporting would miss~\cite{Gneiting2007,Murphy1973}. From a decision-analytic standpoint, this is expected whenever the loss-minimizing action is locally constant in a neighborhood of $p$~\cite{Howard1966,Vickers2006}.

A related observation concerns disagreement priors. Widening uncertain $\pi_i$ changes the prior mass on configurations that include those factors, and therefore changes the myopic VoI ranking in principle. On TEST, however, the inspectable sets are so constrained that ranking changes seldom alter the sequence of realized observations before the budget is exhausted. The small AUC edge for disagreement VoI at $B{=}1$ ($0.873$ versus $0.872$) is consistent with a prior-induced reordering that affects ranking of $p$ more than hard alerts. This is precisely the kind of weak effect that reliability VoI studies report when observation menus are short~\cite{Song2021,Memarzadeh2016,Fauriat2020}. Preserving disagreement remains useful for auditability---analysts can see which factors were treated as contested---even when portfolio-level metrics tie.

Early-warning lead time provides a complementary operational reading. Caught severe projects retain on average ${\approx}3$ remaining schedule updates after the early window, which is enough for managerial escalation in many capital programs but not a guarantee of recovery. Missed severe projects illustrate the opposite boundary: prefix features sometimes lack the cumulative variance signature that the champion relies on, so calibrated probabilities stay below $\tau$. Improving recall further would likely require richer early covariates (contract type, contractor performance histories, or geospatial disruption indicators) that are outside the released NYC schedule tables. That limitation is dataset-bound rather than algorithm-bound~\cite{Tagharobi2025,Egwim2021,Woldesenbet2016}.

Mapping quality sits upstream of both BN parents and several early-warning features. The NLP-only macro-F1 of ${\approx}0.31$ against regex weak labels on VAL would be unacceptable as a sole extractor, which is why the operational pipeline is regex-first. Weighted F1 of $1.0$ on informative regex rows shows that the high-confidence path is stable on the texts it claims. The scientific implication is that free-text delay reasons in open capital data are too formulaic for a purely statistical mapper to invent structure that regex cannot already capture; the BN's value is then less about discovering new linguistic factors and more about coherently combining the factors that regex already identifies under uncertainty~\cite{Onisko2001,Luu2009}.

Calibration and decision curves jointly answer whether the champion's probabilities are usable, not merely rank-correct. ECE $0.030$ and Brier $0.043$ indicate that average predicted severe probability aligns with observed frequencies at the reported resolution, while Fig.~\ref{fig:dca} shows positive net benefit over alert-all across a wide threshold band. In deployment terms, an agency that alerts when $p{\geq}\tau$ is not forced into an extreme operating point to obtain value. That property matters more for triage than a one-point accuracy comparison, even though accuracy remains the community yardstick used in several delay-risk papers~\cite{Vickers2006,GneitingRaftery2007,Khodabakhshian2025,Tagharobi2025}.

Bootstrap intervals quantify sampling uncertainty under the fixed PID split. Accuracy's $95\%$ CI $[0.922,0.958]$ lies entirely above the $0.89$ literature reference, and ROC-AUC's CI $[0.969,0.990]$ remains high. Precision's wider interval $[0.782,0.914]$ reflects minority-class volatility: with only $110$ severe TEST projects, precision is sensitive to small changes in the FP/TP counts. Reporting that width is preferable to claiming a single precision point as stable~\cite{Hand2009,Saito2015}. The PID split histogram confirms that the $70/15/15$ design was realized in the saved artifacts rather than only described in text.

Relative to BBN delay classics, the present evaluation differs in three respects that affect interpretation. First, evidence is observational and multi-period rather than elicited once per project~\cite{Luu2009,Khodabakhshian2025}. Second, the decision problem is budgeted inspection, not only posterior risk reporting~\cite{Memarzadeh2016,Yuan2021}. Third, an independent early-warning model is scored with proper scoring rules and decision curves on the same PIDs~\cite{Gneiting2007,Vickers2006}. These differences mean that numeric gaps versus prior accuracy figures are not pure method rankings; they are joint consequences of data regime and evaluation protocol. The honest scientific claim is therefore comparative under explicit scope, not absolute superiority on a shared benchmark that does not yet exist for open capital-project delay risk.

Failure analysis must also acknowledge computational and modeling boundaries that did not bite here but would bite elsewhere. Exact enumeration at $|\mathcal{F}|{=}8$ hides the cost that would appear if the ontology expanded to dozens of fine-grained causes without Noisy-OR factorization~\cite{Heckerman1996}. Myopic VoI ignores information-gathering value beyond the next inspection, which can matter when inspections are cheap and horizons are long~\cite{Memarzadeh2016,Fauriat2020}. Absence clamping, while faithful to ``never mentioned in reasons,'' can be too aggressive if administrative silence does not imply factor absence. Each of these boundaries suggests a concrete follow-on experiment: larger factor sets with approximate inference, multi-step VoI under explicit inspection costs, and alternative missingness models that treat silence as a latent state rather than a hard zero.

From a portfolio-management perspective, the combined DRIN evidence supports a two-layer workflow. Layer one is the calibrated early-warning score used to rank projects for attention using prefix signals already present in schedule history. Layer two is the BN audit trail used when an analyst can afford a small number of factor inspections and needs an explainable posterior. Layer two should not be sold as a guaranteed policy upgrade over fixed checklists on sparse reason data; it should be sold as a coherent probabilistic account that becomes increasingly useful as inspectable evidence grows. That messaging aligns with the empirical tie in Table~\ref{tab:fullB} and with VoI theory's own warnings about low-information observation menus~\cite{Song2021,Malings2016,Howard1966}.

Finally, the literature gap stated in Section~\ref{sec:lit} can be closed only partially by one metropolitan open-data study. The missing multi-city, multi-agency external validation remains future work. What this study does provide is a reproducible experimental template: PID-held-out splits, regex-auditable factor mapping, Noisy-OR BN with consensus/disagreement priors, budgeted policy scoreboard, and a full probabilistic evaluation pack. Replicating that template on other open capital programs would test whether the VoI-policy tie is a NYC sparsity artifact or a broader property of administrative delay-reason regimes~\cite{Woldesenbet2016,Regona2024,Abioye2021}. Until such replications exist, claims of generalizability remain intentionally restrained.
